
\subsubsection{6.1 Interpretation}

Results document implementation failure, not theoretical refutation. Three detection tiers were specified but not implemented. The core hypothesis—that contamination produces detectable geometric signatures—remains untested.

\subsubsection{6.2 Validation Methodology Observation}

In this research system, data validation (verifying correct dataset loading) occurred independently of algorithm validation (verifying detection logic). This created a testing gap: data correctness was confirmed while processing correctness was not verified.

\textbf{What was checked:} Dataset provenance (real GSM8K and MATH data), training infrastructure (models trained without errors).

\textbf{What was not checked:} Detection algorithm implementation, per-component unit testing, output sensitivity to inputs.

\subsubsection{6.3 Limitations}

\textbf{L1: Complete Implementation Failure (FUNDAMENTAL).} All detection tiers not implemented. This eliminates all research contributions except documenting the failure.

\textbf{L2: Detection Thresholds Never Calibrated (HIGH).} Thresholds were theoretical specifications, not empirically validated on clean baselines. Even correct implementation would require threshold calibration.

\textbf{L3: TSG Paraphrase-Invariance Unvalidated (HIGH).} Core novelty claim never tested. Should be prerequisite validation before full architecture.

\textbf{L4: Geometric Inevitability Untested (HIGH).} Adaptive adversary experiments not executed.

\textbf{L5: Single-Dataset Validation (MEDIUM, acceptable for PoC).} Scoped to GSM8K; generalization unknown but acceptable for existence hypothesis.

\textbf{L6: Statistical Power Insufficient (MEDIUM, correct decision).} N=20 runs executed but detector failure discovered early. Continuing would not fix implementation gaps.

\subsubsection{6.4 Implications}

\textbf{For Contamination Detection Research.} The research gap remains open. Geometric detection methods, TSG paraphrase-invariance, and geometric inevitability are untested hypotheses. Prior work (Fu et al., Dekoninck et al.) remains state-of-the-art.

\textbf{For Validation Methodology.} This case suggests validation strategies verifying data provenance should be complemented by validation verifying processing logic. In this system, per-component unit testing before integration would have detected implementation gaps earlier.
